% TOP_PROBLEM: {"id": 9400055, "title": "Kummer–Vandiver conjecture", "category_id": 1, "category": {"id": 1, "name": "number_theory", "display_name": "Number Theory", "description": "Properties of integers, prime numbers, Diophantine equations.", "slug": "number-theory", "order_index": 1, "created_at": "2026-07-31T15:26:25.670Z"}, "status": "open"}
% FIRST_POSTED: 2026-09-27
% !TeX program = lualatex
% ATTEMPT_STATUS: unresolved
\documentclass[11pt]{article}
\usepackage[margin=25mm]{geometry}
\usepackage{fontspec}
\setmainfont{Latin Modern Roman}
\usepackage{amsmath,amssymb,mathtools}
\usepackage{unicode-math}
\setmathfont{Latin Modern Math}
\usepackage{xurl,hyperref}
\hypersetup{colorlinks=true,urlcolor=blue}
\setcounter{secnumdepth}{0}
\setlength{\parindent}{0pt}
\setlength{\parskip}{5pt}
\setlength{\emergencystretch}{3em}
\begin{document}
\section*{\texorpdfstring{Kummer–Vandiver conjecture}{Kummer–Vandiver conjecture}}\label{9400055}
\subsection{Definitions and mathematical statement}
For a prime $p$, let $\zeta_p=e^{2\pi i/p}$ and let
\[
 K_p^+={\mathbb{Q}}(\zeta_p+\zeta_p^{-1})
\]
be the maximal real subfield of the $p$-th cyclotomic field. Its ideal class group is the group of nonzero fractional ideals of its ring of integers modulo principal fractional ideals; its finite order is the class number $h_p^+$. The conjecture is $p\nmid h_p^+$ for every prime $p$, equivalently that the $p$-primary part of this class group is trivial. The case $p=2$ gives $K_p^+={\mathbb{Q}}$. It does not say $h_p^+=1$ or forbid divisibility by other primes. The catalogue's named 2024 reference is paired with a January 2019 Notices issue URL; that bibliographic mismatch is retained visibly, and Ghate's direct formulation is supplied as an additional source.
\par\smallskip\textit{Formulation and concept references:} \href{https://www.ams.org/journals/notices/201901/201901FullIssue.pdf}{[S1]}, \href{https://mathweb.tifr.res.in/~eghate/vandiver.pdf}{[E1]}.

\subsection{Short English statement}
Does a prime never divide the class number of the maximal real subfield of the cyclotomic field generated by its roots of unity?
\subsection{Research attempt}
\textbf{Status: UNRESOLVED.} This attempt does not prove that the plus class group has trivial $p$-primary part for every prime. It proves small-prime cases by explicit Minkowski bounds, explains precisely how the plus part sits inside the full cyclotomic class group, and separates the Bernoulli-number irregularity test from the Vandiver condition. None of the finite calculations is represented as a proof for arbitrary primes.

\paragraph{Source audit and notation.}
The inherited bibliography already records the mismatch between the title of [S1] and its linked Notices issue. That mismatch is preserved rather than silently replaced. Ghate's author-hosted notes [E1], checked on 27 September 2026, directly state Vandiver's conjecture and distinguish its even-character components from the odd components governed by Bernoulli-number criteria. Their historical numerical bound is not used as a current computation record. For the Minkowski method, I also use Conrad's author-hosted notes [E2].

For odd $p$, write $K=\mathbb Q(\zeta_p)$, $F=K_p^+$, $d=[F:\mathbb Q]=(p-1)/2$, and let $\iota$ denote complex conjugation on $K$. Class groups are written multiplicatively except when additive notation simplifies their $p$-primary decomposition. The class number $h_p^+$ may be greater than one while still satisfying the conjecture. Conversely, finding an ideal which is not principal does not refute Vandiver unless its class contributes nontrivial $p$-torsion.

\paragraph{1. The real cyclotomic discriminant and a finite class-group criterion.}
I use the standard cyclotomic integral-basis facts
\[
 \mathcal O_K=\mathbb Z[\zeta_p],\qquad
 \mathcal O_F=\mathbb Z[\zeta_p+\zeta_p^{-1}],\qquad
 |D_K|=p^{p-2}.
\]
The relation $\zeta_p^2-(\zeta_p+\zeta_p^{-1})\zeta_p+1=0$ makes $\mathcal O_K$ free over $\mathcal O_F$ with basis $1,\zeta_p$. Its relative discriminant is generated by $(\zeta_p-\zeta_p^{-1})^2$. The absolute norm of that relative discriminant is $p$: indeed, multiplying over one representative from each pair $a,-a$ gives the same absolute product as
\[
 \prod_{a=1}^{p-1}(1-\zeta_p^a)=p,
\]
using that multiplication by two permutes nonzero residues modulo $p$. The discriminant tower formula therefore gives
\[
 |D_K|=D_F^2p,\qquad D_F=p^{(p-3)/2}.
\]
The integral-basis and discriminant-tower theorems are classical algebraic-number-theory inputs; the calculation identifies their specific consequence needed here.

Minkowski's ideal-class theorem says that every ideal class in a totally real field of degree $d$ and discriminant $D_F$ contains an integral ideal of norm at most
\[
 B_p=\frac{d!}{d^d}\sqrt{D_F}.
\]
If $B_p<2$, the only possible norm is one, and the ideal is the entire ring, so every class is principal. More generally, if no prime ideal has norm between 2 and $B_p$, every integral ideal with norm at most $B_p$ must be the unit ideal, by unique factorization of ideals and multiplicativity of norm. This again proves class number one. Both criteria are stronger than the requested nondivisibility, but useful in small degrees.

\paragraph{2. Explicit proofs for $p\le13$.}
For $p=2$ and $p=3$, the real subfield is $\mathbb Q$, whose ideal class group is trivial. For $p=5$ and $p=7$, the bounds are
\[
 B_5=\frac{\sqrt5}{2}<2,
 \qquad B_7=\frac{14}{9}<2.
\]
Thus $h_5^+=h_7^+=1$.

For the next cases use the splitting law for an unramified rational prime $q\ne p$. The Galois group of $F/\mathbb Q$ is $(\mathbb Z/p\mathbb Z)^\times/\{\pm1\}$, and Frobenius at $q$ is the class of $q$. Consequently every prime ideal over $q$ has residue degree
\[
 f_p(q)=\min\{f\ge1:q^f\equiv1\text{ or }-1\pmod p\},
\]
and norm $q^{f_p(q)}$. This is the standard Frobenius splitting theorem applied to this explicit quotient group, not a claim that residue degree equals the full multiplicative order modulo $p$.

At $p=11$, one has $d=5$, $D_F=11^4$, and
\[
 B_{11}=\frac{120\cdot121}{5^5}=\frac{14520}{3125}<5.
\]
Only the rational primes 2 and 3 can divide the norm of a nontrivial ideal with norm at most this bound. Both have $f_{11}(q)=5$: their first powers up to exponent four are not $\pm1$ modulo 11, while the fifth is $\pm1$. Their prime-ideal norms are therefore $2^5=32$ and $3^5=243$, both too large. The ramified prime 11 is also above the bound. Hence every ideal class has a representative of norm one, and $h_{11}^+=1$.

At $p=13$, one has $d=6$, $D_F=13^5$, and
\[
 B_{13}=\frac{720}{6^6}\,13^{5/2}<10.
\]
The inequality can be verified by squaring positive quantities and using integers; a decimal approximation is unnecessary. The only candidate rational primes are $2,3,5,7$, with degrees
\[
 \bigl(f_{13}(2),f_{13}(3),f_{13}(5),f_{13}(7)\bigr)
   =(6,3,2,6).
\]
Their prime-ideal norms are $64,27,25,117649$, respectively, all greater than 10. The ramified prime 13 cannot occur below the bound either. Thus $h_{13}^+=1$. These computations prove Vandiver for $p=2,3,5,7,11,13$, with every needed small-prime exclusion explicit.

The method need not remain this simple as $p$ grows. Once small prime ideals occur within the Minkowski bound, one must prove their classes principal or understand their relations. Merely listing the ideals is not a class-group computation, and the rapidly growing discriminant prevents the same short argument from becoming a uniform proof.

\paragraph{3. Extension and norm identify the plus $p$-primary part.}
Let $j:\operatorname{Cl}(F)\to\operatorname{Cl}(K)$ extend ideals and let $N:\operatorname{Cl}(K)\to\operatorname{Cl}(F)$ be the ideal norm. These are well-defined on classes because principal ideals extend and norm to principal ideals. Since $[K:F]=2$,
\[
 N(j(c))=c^2.
\]
For odd $p$, squaring is an automorphism of every finite abelian $p$-group. It follows immediately that $j$ is injective on the $p$-primary subgroup: if $j(c)=1$ and $c$ has $p$-power order, then $c^2=1$, forcing $c=1$.

Write $A$ for the $p$-primary subgroup of $\operatorname{Cl}(K)$ and
\[
 A^+=\{a\in A:\iota(a)=a\}.
\]
The image under $j$ of a class from $F$ is conjugation-invariant. Conversely, for a class $a\in A$ one has the ideal-class identity
\[
 j(N(a))=a\,\iota(a).
\]
For $a\in A^+$, its right side is $a^2$. The unique square root operation on the finite $p$-group then shows that every member of $A^+$ lies in the image of the $p$-primary subgroup of $\operatorname{Cl}(F)$. More explicitly, choose an integer $r$ with $2r\equiv1\pmod{p^e}$, where $p^e$ kills $A$. Then $a=j(N(a)^r)$. Consequently
\[
 \operatorname{Cl}(F)[p^\infty]\cong A^+.
\]
This gives the precise full-field formulation of Vandiver: the plus $p$-primary subgroup is zero.

In additive notation, division by two is legitimate on $A$, so every $a\in A$ decomposes uniquely as
\[
 a=\frac{a+\iota(a)}2+\frac{a-\iota(a)}2,
 \qquad A=A^+\oplus A^-.
\]
The two summands are the $+1$ and $-1$ eigenspaces of conjugation. Triviality of all of $A$ implies Vandiver, but nontriviality of $A^-$ does not by itself imply nontriviality of $A^+$. Confusing these two implications is a common route to a spurious counterexample.

\paragraph{4. What Bernoulli numbers do and do not test.}
For a prime $p$, classical regularity criteria involve divisibility by $p$ of numerators of the Bernoulli numbers $B_2,B_4,\ldots,B_{p-3}$. In the more refined character decomposition described in [E1], the relevant Bernoulli conditions detect odd-character components. Vandiver concerns the even components. The two questions are related by deep reflection results, but they are not the same test with a different name.

For a concrete warning, the exact rational value
\[
 B_{32}=-\frac{7709321041217}{510}
\]
has numerator divisible by 37, while its denominator is not. This makes 37 an irregular prime. It does not, just from this divisibility, establish $37\mid h_{37}^+$. Treating every irregular prime as a Vandiver counterexample would mistakenly transfer information about the wrong class-group part.

The distinction is already visible in abstract finite-group algebra: a group of order $p$ on which an involution acts by inversion has nontrivial minus part and trivial plus part. This example is not asserted to realize an arbitrary cyclotomic class group; it simply shows that nontriviality of a group with involution does not logically force a nontrivial fixed subgroup. Arithmetic information beyond total $p$-divisibility is indispensable.

\paragraph{5. Why a naive norm descent does not finish the problem.}
Suppose $c\in\operatorname{Cl}(F)$ has order $p$. Choose an ideal $\mathfrak a$ representing it. There is then an algebraic integer or fractional generator $\alpha$ with $\mathfrak a^p=(\alpha)$, after the usual choice of a fractional ideal representative. This identity does not imply that $\mathfrak a$ is principal: extracting a $p$th root of a principal ideal generator inside $F$ is exactly the issue. Taking norms gives a numerical perfect-power relation which forgets the ideal class, rather than making it vanish.

Likewise, the equality $N\circ j=[2]$ proves injectivity on the $p$-part, not its triviality. To deduce triviality one would need an independent theorem that the invariant $p$-classes upstairs vanish. Substituting this desired conclusion as a property of the norm map is circular. Cyclotomic units and character decompositions provide much finer arithmetic structure, but no uniform argument annihilating all the even components is supplied here.

\paragraph{6. Exact checks and remaining gap.}
The accompanying arithmetic diagnostic evaluates the stated discriminant and Minkowski inequalities without rounding, computes residue degrees by modular exponentiation, and checks every possible rational prime below the displayed bounds. It also computes Bernoulli numbers by their rational recurrence and verifies the numerator divisibility at $(p,2m)=(37,32)$. Finite abelian examples check the plus/minus and squaring identities. The script is not a general class-group algorithm and does not claim a new numerical verification range for Vandiver.

The proved small-prime cases and the extension-norm equivalence sharpen the target but leave its central arithmetic content untouched: rule out an element of order $p$ in the plus part for every odd prime $p$. Neither a rapidly expanding Minkowski search nor a test of the irregular Bernoulli indices establishes that statement. The unrestricted conjecture therefore remains unresolved by this attempt.

\subsection{Sources}
\begin{itemize}
\item[S1]Hou Rong Qin, The Vandiver Conjecture and a New Conjecture on the Distribution of Irregular Primes, Acta Mathematica Sinica, Chinese Series 67(2) (2024), 341-346.
\url{https://www.ams.org/journals/notices/201901/201901FullIssue.pdf}.
\textit{Formulation source.}
\item[E1]Eknath Ghate, Vandiver’s Conjecture via K-theory, Summer School on Cyclotomic Fields, Pune, 7–30 June 1999, Conjecture 1.
\url{https://mathweb.tifr.res.in/~eghate/vandiver.pdf}.
\textit{Added primary source: Direct statement for the class number of the maximal real cyclotomic subfield; added because the catalogue citation and URL do not match.. Consulted 24 September 2026.}
\item[E2]Keith Conrad, Class group calculations by Minkowski bound, author-hosted notes.
\url{https://kconrad.math.uconn.edu/blurbs/gradnumthy/classgpex.pdf}.
\textit{Added primary expository source: Ideal-class representatives of bounded norm and the method of checking prime ideals. Consulted 27 September 2026.}
\end{itemize}
\end{document}
